The primary expected outcome is a validated LPV-LFR grey-box augmented model that improves settling prediction accuracy over the baseline, particularly for cross-coupling and position-dependent dynamics, with BFR reported on held-out operating points and motion profiles in both simulation and experimental validation. Results will also clarify whether an LPV baseline outperforms frozen LTI alternatives, and whether orthogonal projection-based regularization preserves physical parameter interpretability in the LPV-LFR setting, with implications for extending orthogonal projection regularization to coupled MIMO and LFR settings. The resulting methodology provides ASMPT with a structured augmentation approach that can be evaluated against current practice and considered for further development or implementation. Beyond the graduation report, deliverables include the identified models, implementation code, validation scripts, and a documented modeling workflow.